\documentclass[10pt,a4paper]{amsart}
\usepackage[margin=19mm]{geometry}
\usepackage{amsmath,amssymb,mathtools}
\usepackage[scr=rsfs,cal=euler]{mathalfa}
\usepackage{xfrac}
\usepackage[unicode,hidelinks,bookmarks=false]{hyperref}
\newcommand{\ie}{i.e.\ }
\newcommand{\cf}{cf.\ }
\newcommand{\resp}{respectively\ }
\newcommand{\Subsection}{\subsection}
\providecommand{\nobreakdash}{\nobreak\hbox{-}}
\setlength{\emergencystretch}{2em}
\multlinegap0pt
\usepackage{bm}
\usepackage{bbm}
\usepackage{dsfont}
\usepackage{xspace}
\usepackage{cancel}
\usepackage{mathtools}
\usepackage{amsthm}
\makeatletter
\@ifundefined{theorem}{\newtheorem{theorem}{Theorem}}{}
\@ifundefined{definition}{\newtheorem{definition}[theorem]{Definition}}{}
\@ifundefined{lemma}{\newtheorem{lemma}[theorem]{Lemma}}{}
\makeatother
\title[An iterative construction of complete K\"ahler--Einstein met]{An iterative construction of complete K\"ahler--Einstein metrics}
\author{Quang-Tuan Dang, Tat Dat T\^o}
\date{Source: arXiv:2410.12599v3 (CC BY 4.0)}
\begin{document}
\maketitle

\noindent\emph{Source and licence.} An iterative construction of complete K\"ahler--Einstein metrics. Version arXiv:2410.12599v3 (CC BY 4.0), distributed under the Creative Commons Attribution 4.0 International licence, as recorded in the arXiv licence metadata for this exact version and shown on the abstract page https://arxiv.org/abs/2410.12599v3. Full rights notice: licenses/arxiv-2410.12599-CC-BY-4.0.txt.\par\smallskip

\noindent\textit{Editorial scope.} Counted unit: the Lemma whose source label is \texttt{lem:normal\_chart} in the article (source0.tex lines 442--445, source label \texttt{lem:normal\_chart}), delivered together with the article's own complete original proof of it (source0.tex lines 446--462, one \texttt{proof} environment of 17 lines). No mathematics is written, repaired or re-derived: statement and proof are the article's own text line for line. Required context retained uncounted: def:bd\_geo (lines 427--439). Every definition, display and label of those lines is retained unchanged. Omitted from this package and not redistributed: 1--426, 440--441, 463--1257. Nothing inside a retained excerpt is omitted or altered: every retained body is delivered exactly as the article's lines are written, so no omission receipt of any kind accompanies this package. Citations: the retained excerpts carry no citation command at all, so no bibliography entry of the article is transcribed and no reference is redistributed. Reference adaptation: none was needed and none was made; every reference inside the delivered unit resolves inside it, so no unresolved reference or citation is produced. Printed slips kept literal: no printed word, symbol, hypothesis or number of the counted statement or of its proof was altered. Layout and class: the article's own class is \texttt{amsart} with packages including inputenc, babel, fancyhdr, amsmath, amsfonts, amstext, mathrsfs, amssymb, amsthm, amscd, amsxtra, wasysym, graphicx, esint; the wrapper is the lane's pinned \texttt{amsart} editorial wrapper at 10pt on A4, which restates the theorem-like environments the retained text needs and adds the article's own macro definitions that the retained text needs (none), with extra wrapper packages (bm, bbm, dsfont, xspace, cancel, mathtools). The delivered numbering is the unit's own. Assets: no retained span contains a figure environment, a graphics-inclusion command, a PDF-page inclusion, a table or a TikZ picture; no image file is copied into this package, the package declares no asset dependency and no image file is redistributed with it. Licence: version arXiv:2410.12599v3 of this article is distributed under the Creative Commons Attribution 4.0 International licence, as recorded in the arXiv licence metadata for this exact version and shown on the abstract page https://arxiv.org/abs/2410.12599v3. A full rights notice is delivered at licenses/arxiv-2410.12599-CC-BY-4.0.txt. Characters: every retained excerpt is pure ASCII, so the delivered engine, XeLaTeX with its default Latin Modern text font, reports no missing character.
\begin{definition}\label{def:bd_geo}
Suppose $(X,\omega)$ is a $n$-dimensional K\"ahler manifold. We say that $(X,\omega)$ has {\em bounded geometry} of order $k\in\mathbb{N}$ if there exists $0<r$ such that for any $p\in X$ there is  a domain $U$ in $\mathbb{C}^n$ and a  holomorphic map   $\psi:U  \rightarrow X $ of maximal rank everywhere,  satisfying 
\begin{enumerate}
    \item $B_{\mathbb{C}^n}(0,r)\subset U $ and $\psi(0)=p;$
    
    \item  On $U$, there exists a constant $c>0$ depending only on $n$, $r$ such that $$c^{-1}\omega_{\mathbb{C}^n}\leq \psi^*\omega \leq c\omega_{\mathbb{C}^n };$$
    \item  there is a constant $A_k>0$ depending only on $k$ such that
$$\sup_{x\in B_{\mathbb{C}^n}(0,r)}\left| \frac{\partial^{|\alpha|+|\beta|}g_{i\bar j}(x)}{\partial z^\alpha \partial \bar z^\beta}\right|\leq A_k, \;\;\forall\;|\alpha|+|\beta|\leq k,$$
\end{enumerate}
where $g_{i\bar j}$ is the component of $\psi^*\omega$ on $U$ in terms of natural coordinates $(z^1, \ldots ,z^n)$, and $\alpha$, $\beta$ are the multiple indices with $|\alpha|=\alpha_1+\cdots+\alpha_n$.

The map $\psi$ is called a {\em quasi-coordinate map} and the pair $(U,\psi)$ is called a {\em quasi-coordinate chart} of $X$. % The number $r$ is called {\em radius of quasi-bounded geometry}.
\end{definition}

\begin{lemma}\label{lem:normal_chart}
Suppose $(X,\omega)$ is a K\"ahler manifold with bounded geometry of order at least $ 5$.  Then there exists   $c>0$, $A_k>0$, for $k=1,2, \ldots $  such that for  any point $p\in X$,  there exist $0<\varepsilon=\varepsilon(p)< r$,  and a holomorphic map  $\kappa:  B_{\mathbb{C}^n}(0,\varepsilon) \rightarrow X $ which is biholomorphic on its image  with $\kappa(0)=p$, $\kappa^* \omega=dd^c \varphi$ satisfying the conditions (2), (3) in  Definition~\ref{def:bd_geo}, and $\varphi_{a\overline{b}}(0)=\delta_{a b}$, 
$\varphi_{a\overline{b} c}(0)=\varphi_{a\overline{b}\overline{c}}(0)=0$, and  $\varphi_{a \overline{b}\overline{c} \overline{d}}(0)=\varphi_{a b c\overline {d}}(0)=0$. 
\end{lemma}

\begin{proof} 
For any $p\in X$, we fix a local quasi-coordinate chart $(U,\psi)$ at $p$  as in Definition~\ref{def:bd_geo}.  By complex linear transformation we can assume that $\psi^*\omega =dd^c \phi$ on $U$ with $\phi_{i \bar j}(0)=\delta_{ij}$.  Since $\psi$ is of maximal rank everywhere, there exists $0<\varepsilon_1=\varepsilon_1(p)<r$ such that $\psi: B(0, \varepsilon_1)\rightarrow 
\psi(B(0, \varepsilon_1))\subset X$ is biholomorphic.  We can choose $\varepsilon_2<\varepsilon_1$ depending on $A^i_{k\ell}, B^i_{mnp} $  such that the holomorphic map  $\Tilde{\psi}: B(0,\varepsilon_2) \rightarrow B(0,\varepsilon_1) $,  $z^i= \Tilde\psi(w)^i= w^i+ A^i_{k\ell}w^kw^\ell + B^i_{mnp} w^m w^n w^p$   is biholomorphic on its image, where  the coefficients $A^i_{k\ell},B^i_{mnp} $  will be chosen hereafter. 

\medskip
Now we have  $\kappa=\psi\circ \tilde\psi:B(0,\varepsilon_2)\rightarrow X$ is biholomorphic on its image,  and the pull-back metric $ \kappa^*\omega =dd^c \varphi$  satisfies
\begin{align*}
   \varphi_{a\bar b} (w)&=\phi_{a\bar b}(\tilde \psi(w)) + 2A^i_{a k}\phi_{i\bar b}w^{k}  + 2\overline{A^j_{b\ell}\phi_{a\bar j}}\overline{w}^\ell + 3B^i_{amn}w^m w^n  \phi_{i\bar b} + 3 \overline{B^j_{bpq}}\phi_{a\bar j}\overline{w}^p\overline{w}^q \\
   &\quad + 4 A^{i}_{a k}\overline{A^j_{b\ell}}\phi_{i\bar j}w^k\overline{w}^\ell+O(|w|^3). 
\end{align*}
Therefore we have
$$
\varphi_{a\bar b c}(0) = \phi_{a\bar b c}(0)+ 2 A^{b}_{ac};  \varphi_{a\bar b \bar c}(0) = \phi_{a\bar b \bar c}+ 2 \overline{ A^{b}_{ac}}; \varphi_{a\bar b cd}(0) = \phi_{a\bar b c d}(0) +6 B^{b}_{acd}+C,
$$
where $C$ only depends on $A^{b}_{ac}$ and $\phi_{a\bar b c}(0)$.
By choosing $A^b_{ac}=-\frac{1}{2}\phi_{a\bar b c}(0)$ and then $B^b_{acd}=-\frac{1}{6}\phi_{a\bar b c d}(0) - C$,  we obtain the local normal coordinates for any $p\in X$. In particular, $A^{a}_{bc}$ and $B^a_{bcd}$ are uniformly bounded for any $p\in X$ by the  definition of bounded geometry. Since $\varepsilon_2<\varepsilon_1$ depends only on $A^{a}_{bc}$ and $B^a_{bcd}$,  we get the map $\kappa: B(0,\varepsilon_2)\rightarrow X$ as desired.
\end{proof}

\end{document}
